\chapter{综合方案推荐与结论}\label{chap:recommendation}

\section{综合选型原则}

信道编码方案应首先按业务类型分流。200 bit与300 bit低速电文采用BCH候选，300 bit高速电文采用卷积码或LDPC候选；三类编码承载的业务、发送长度和处理机制不同，不建立跨业务统一排名。在业务边界确定后，系统依次检查目标FER覆盖、实际发送长度与码率、实时处理和实现资源、信道适应性，最后决定是否配置交织。该顺序延续表\ref{tab:bch-conclusions}、表\ref{tab:cc-main-conclusions}、表\ref{tab:cc-system-recommendations}、表\ref{tab:ldpc-length-tradeoff}和表\ref{tab:ldpc-decoder-tradeoff}已经冻结的码族内部结论。

可靠性是进入后续比较的前置条件。方案若在规定信噪比范围内不能覆盖目标FER，或在目标场景中形成明显错误平台，则不能仅凭较高码率或较短软件CPU时间成为主推荐。可靠性满足后，再比较payload之外的tail、filler、shortening和puncturing所形成的实际发送长度。实时性评价同时考虑首次输出、P95、最大观测软件时间、迭代上限、回溯深度和结构缓存；其中CPU时间只描述当前处理器、编译器和计时边界，不能直接等同于硬件时延。

信道适应性依据第7章的AWGN、固定多径、固定频偏、线性时变频偏、已知短时遮挡和未知突发结果，近码率横向选择以表\ref{tab:s5-recommendations}为最高优先级证据。交织不是常驻模块，只在连续错误集中出现且已有正式试验支持时启用。表\ref{tab:s7-engineering-recommendation}所支持的BCH与卷积码配置均受第9章具体码型和突发比例约束，LDPC则保持无交织。

\begin{table}[htbp]
  \centering\scriptsize
  \caption{综合选型判据与证据来源}
  \label{tab:final-selection-criteria}
  \resizebox{\textwidth}{!}{\begin{tabular}{c l l l l}
    \toprule
    决策层级 & 工程问题 & 核心指标 & 主要证据章节 & 使用边界 \\
    \midrule
    业务 & 低速或高速、payload长度 & 200/300 bit、业务接口 & 第2、4--7章 & 禁止跨业务统一排名 \\
    可靠性 & 是否覆盖目标工作点 & FER、BER、错误平台 & 表\ref{tab:bch-conclusions}、\ref{tab:s5-recommendations} & 未覆盖目标者不进入效率优选 \\
    传输效率 & 每帧实际占用多少符号 & 发送长度、实际码率、填充与打孔 & 表\ref{tab:cc-main-conclusions}、\ref{tab:ldpc-length-tradeoff} & 固定$E_s/N_0$与固定$E_b/N_0$不可混用 \\
    实时性 & 何时产生稳定输出 & 首输出、P95、最大观测、迭代与回溯 & 表\ref{tab:cc-system-recommendations}、\ref{tab:s6-engineering-summary} & 软件CPU时间不等于物理链路时延 \\
    实现资源 & 译码与缓存是否可承受 & 状态数、图规模、存储、buffer & 表\ref{tab:s6-engineering-summary} & 不统一折算异构操作次数 \\
    信道适应性 & 非理想信道下是否保持覆盖 & 六类信道FER门限与损失 & 表\ref{tab:s5-recommendations} & 仅适用于第7章规定模型 \\
    突发与交织 & 是否需要重排错误位置 & 平均/最坏FER、跨度、缓存 & 表\ref{tab:s7-engineering-recommendation} & 仅外推到已验证码型与突发比例 \\
    \bottomrule
  \end{tabular}}
\end{table}

\section{低速电文BCH综合推荐}

200 bit业务形成低复杂度、高可靠性和明显连续突发三条路线。资源受限且信道以分散错误为主时，主推荐$19\times\mathrm{BCH}(15,11,1)$：200 bit payload补9 bit后形成19个11 bit子块，发送285 bit，实际码率为$200/285\approx0.701754$，采用综合校验查表译码。该路线结构短、查表路径直接，在当前软件实现中译码时间较低；代价是AWGN可靠性低于较长整块BCH。

常规信道可靠性优先时，主推荐缩短$\mathrm{BCH}(511,385)$，发送326 bit、$t=14$，实际码率为$200/326\approx0.613497$，采用BM与Chien搜索。该配置以更多发送冗余和较高代数译码开销换取第4章正式AWGN结果中的较高可靠性。对明显连续突发，结论只适用于上述200 bit分块方案：性能优先采用行列交织R15，结构组织优先可采用D19；两者均缓存完整285 coded bits。第9章并未验证200 bit整块BCH或300 bit整块BCH的同类交织，因此不得据此扩展。

300 bit业务不新增交织路线。综合平衡主推荐缩短$\mathrm{BCH}(511,421)$，发送390 bit、$t=10$，实际码率$300/390\approx0.769231$；可靠性优先采用缩短$\mathrm{BCH}(511,385)$，发送426 bit、$t=14$，实际码率$300/426\approx0.704225$。两者均采用BM与Chien搜索，选择实质是90 bit与126 bit冗余之间的可靠性和发送占用取舍。

\begin{table}[htbp]
  \centering\scriptsize
  \caption{低速BCH最终推荐矩阵}
  \label{tab:final-bch-matrix}
  \resizebox{\textwidth}{!}{\begin{tabular}{c l l l c c l l}
    \toprule
    信息长度 & 目标/场景 & 主推荐 & 译码器/交织 & 发送长度 & 实际码率 & 主要依据 & 主要代价或边界 \\
    \midrule
    200 & 低复杂度 & $19\times\mathrm{BCH}(15,11,1)$ & 查表/NONE & 285 & 0.701754 & 短码与最低软件译码时间 & AWGN可靠性低于长整块方案 \\
    200 & 常规可靠性 & 缩短$\mathrm{BCH}(511,385)$，$t=14$ & BM+Chien/NONE & 326 & 0.613497 & 第4章AWGN可靠性优先 & 冗余与代数译码开销较高 \\
    200 & 明显连续突发 & 分块BCH & 查表/R15，D19备选 & 285 & 0.701754 & 表\ref{tab:s7-engineering-recommendation} & 仅限200 bit分块方案，缓存285 bit \\
    300 & 综合平衡 & 缩短$\mathrm{BCH}(511,421)$，$t=10$ & BM+Chien/NONE & 390 & 0.769231 & 表\ref{tab:bch-conclusions} & 不含正式交织验证 \\
    300 & 可靠性优先 & 缩短$\mathrm{BCH}(511,385)$，$t=14$ & BM+Chien/NONE & 426 & 0.704225 & 更强纠错能力 & 发送占用与译码开销增加 \\
    \bottomrule
  \end{tabular}}
\end{table}

\section{高速电文卷积码综合推荐}

高速卷积码统一采用300 bit payload、约束长度$K=7$、八进制生成多项式171/133、6 bit零尾和64个状态。软判决Viterbi为性能主路线，硬判决仅作为接收接口或资源极紧条件下的备选。表\ref{tab:cc-main-conclusions}显示，在FER$=0.1$处，浮点软判决相对硬判决的$E_s/N_0$收益为1.857--2.085 dB；这一收益来自更完整的幅度可靠性，而不增加发送冗余。

五类系统配置服务于不同基准，不构成统一排名。可靠性优先采用$1/2$母码、612 bit发送和整块浮点软判决；吞吐优先采用约$3/4$打孔、408 bit发送及$W=160,S=25,D=126$；时延优先采用约$2/3$打孔、459 bit发送及$W=128,S=25,D=84$；存储优先和综合平衡均采用$1/2$母码与$W=96,S=16,D=70$。高码率配置只有在链路FER余量满足时才具有工程意义。

连续时隙需要分批交付时，优先采用$50\times6$组织。$50\times6$、$100\times3$和$150\times2$产生相同最终码流，FER与BER不因组织方式改变；选择$50\times6$依据的是等待分布和P95输出时间，而不是更高可靠性。定点软信息实现优先采用Q8，第5章在FER$=0.1$处给出的相对浮点损失处于约0.008 dB以内量级，但系统性能基准仍以浮点软判决为准。

\begin{table}[htbp]
  \centering\scriptsize
  \caption{卷积码系统级最终配置矩阵}
  \label{tab:final-cc-matrix}
  \resizebox{\textwidth}{!}{\begin{tabular}{l c c l c c c l l}
    \toprule
    工程目标 & 实际码率 & 发送长度 & 译码方式 & $W$ & $S$ & $D$ & 时隙/软信息 & 选择理由与边界 \\
    \midrule
    可靠性优先 & 0.4902 & 612 & 整块浮点软Viterbi & -- & -- & 306 & 整帧/浮点 & 冗余充分；等待完整码字 \\
    吞吐优先 & 0.7353 & 408 & 滑窗软Viterbi & 160 & 25 & 126 & 连续/浮点 & 发送最短；仅在FER余量允许时使用 \\
    时延优先 & 0.6536 & 459 & 滑窗软Viterbi & 128 & 25 & 84 & $50\times6$/浮点 & 较早分批输出；可靠性低于$1/2$母码 \\
    存储优先 & 0.4902 & 612 & 滑窗软Viterbi & 96 & 16 & 70 & 连续/浮点 & 覆盖既定FER门限且译码存储较低 \\
    综合平衡 & 0.4902 & 612 & 滑窗软Viterbi & 96 & 16 & 70 & $50\times6$/Q8优先 & Q8为实现配置，浮点仍是性能基准 \\
    \bottomrule
  \end{tabular}}
\end{table}

\section{高速电文LDPC综合推荐}

三种BG2 Direct QC-LDPC短块配置均承载300 bit payload，最大迭代次数为32。N480发送480 bit、实际码率0.6250，适用于传输效率优先；N560发送560 bit、实际码率0.5357，其价值是接近576 bit目标的接口适配，不作为默认综合最优；N640发送640 bit、实际码率0.46875，适用于可靠性和弱信号优先。归一化系数必须随码长配置：N480和N560取$\alpha=0.95$，N640取$\alpha=0.80$。

工程译码优先采用layered NMS，BP保留为性能参考。在第8章统一的N560对比中，NMS相对BP的目标FER性能损失小于约0.04 dB，同时当前软件处理时间显著下降；更广的第6章结果亦表明两者目标FER门限接近。该结论只说明当前算法、数据结构和软件环境下NMS具有更合适的实现取舍，不表示其在所有平台或所有迭代调度下全面优于BP。

\begin{table}[htbp]
  \centering\scriptsize
  \caption{LDPC最终推荐配置}
  \label{tab:final-ldpc-matrix}
  \resizebox{\textwidth}{!}{\begin{tabular}{l c c l c c l l}
    \toprule
    目标 & 码长 & 实际码率 & 译码器 & $\alpha$ & 最大迭代 & 主要优势 & 主要代价与定位 \\
    \midrule
    传输效率 & 480 & 0.6250 & layered NMS & 0.95 & 32 & LDPC候选中发送最短 & 可靠性低于N640 \\
    接口适配 & 560 & 0.5357 & layered NMS & 0.95 & 32 & 接近576 bit接口目标 & 非默认性能折中 \\
    可靠性/弱信号 & 640 & 0.46875 & layered NMS & 0.80 & 32 & 更完整的低FER覆盖 & 发送与校验图规模较大 \\
    性能参考 & 480/560/640 & 对应码率 & layered BP & -- & 32 & 和积算法基准 & 当前软件时间明显高于NMS \\
    \bottomrule
  \end{tabular}}
\end{table}

\section{高速电文卷积码与LDPC场景化选择}

近$2/3$组比较CC R2/3（实际码率约0.6536）与LDPC N480（0.6250），属于近码率而非完全同码率比较。表\ref{tab:s5-recommendations}表明，在AWGN、固定多径、固定频偏和线性时变频偏下，CC R2/3的绝对FER门限均更低，因而成为常规全帧信道主推荐；在已知5\%短时遮挡下，N480保持FER$=0.1$目标覆盖，而CC未形成相同覆盖，N480转为该场景的主候选。未知5\%突发干扰下，两者均未形成稳定低FER结果，当前无已验证的稳定低FER主方案。

近$1/2$组比较CC R1/2（约0.4902）与LDPC N640（0.46875）。N640在AWGN、固定多径、固定频偏和线性时变频偏下具有更低绝对FER门限，并在已知5\%遮挡下形成更完整的目标覆盖，可靠性和弱信号优先时应采用N640。CC R1/2仍具有固定64状态、无需迭代、可滑窗和可分批输出等实时实现优势，适用于连续处理或可预测等待优先的系统。未知5\%突发下，N640只表现出有限FER$=0.1$覆盖且明显退化，CC未形成相同低FER覆盖，因此不把任一方案写成稳定主推荐。

\begin{table}[htbp]
  \centering\scriptsize
  \caption{高速电文场景化方案选择矩阵}
  \label{tab:final-highspeed-scenarios}
  \resizebox{\textwidth}{!}{\begin{tabular}{c l l l l l}
    \toprule
    目标码率组 & 信道场景 & 主推荐 & 备选 & 可靠性与实现依据 & 边界 \\
    \midrule
    近$2/3$ & AWGN/固定多径 & CC R2/3软Viterbi & N480 NMS & CC绝对FER门限更低，固定64状态 & 近码率比较 \\
    近$2/3$ & 固定/时变频偏 & CC R2/3软Viterbi & N480 NMS & CC绝对FER门限更低 & 依赖规定帧内相位轨迹 \\
    近$2/3$ & 已知5\%遮挡 & N480 NMS & -- & N480保持FER$=0.1$覆盖 & 受损位置已知 \\
    近$2/3$ & 未知5\%突发 & 无已验证稳定低FER主方案 & -- & 两者均未形成稳定低FER结果 & 不能强行外推 \\
    \midrule
    近$1/2$ & AWGN/固定多径 & N640 NMS & CC R1/2软Viterbi & N640绝对FER门限更低 & CC适合实时连续处理 \\
    近$1/2$ & 固定/时变频偏 & N640 NMS & CC R1/2软Viterbi & N640绝对FER门限更低 & 依赖规定帧内相位轨迹 \\
    近$1/2$ & 已知5\%遮挡 & N640 NMS & -- & N640目标覆盖更完整 & 受损位置已知 \\
    近$1/2$ & 未知5\%突发 & 无已验证稳定低FER主方案 & N640性能参考 & N640仅有限FER$=0.1$覆盖 & 两者均明显退化 \\
    \bottomrule
  \end{tabular}}
\end{table}

\section{译码、连续处理与交织联合配置}

译码器与交织器必须同具体码型和业务条件联合选择。低复杂度200 bit分块BCH在分散错误下使用查表且不交织，明显连续突发时配置R15，D19作为结构清晰的备选；整块BCH采用BM与Chien搜索且不交织。卷积码常规链路采用软Viterbi并关闭交织，需要连续输出时按目标选取表\ref{tab:final-cc-matrix}中的滑窗和时隙参数。

卷积码在2\%突发下可按资源切换：缓存优先采用D8，跨度64 trellis steps、缓存128 coded bits；FER优先采用Pseudo128，跨度128 trellis steps、缓存256 coded bits。D8与Pseudo128是不同资源规模的工程配置比较；只有D16与Pseudo128构成相同128 trellis-step跨度和256 coded-bit缓存下的受控排列比较。5\%突发下全部正式卷积码交织配置仍接近FER饱和，当前未形成满足低FER要求的稳定配置。LDPC始终采用NONE interleaver与NMS；第9章LDPC数据只承担无突发AWGN近码率基线，不能据此推断其突发性能。

\begin{table}[htbp]
  \centering\scriptsize
  \caption{译码与交织联合配置}
  \label{tab:final-decoder-interleaver}
  \resizebox{\textwidth}{!}{\begin{tabular}{l l l l c l l}
    \toprule
    编码方案 & 常规译码 & 连续处理 & 突发配置 & 结构缓存/bit & 适用条件 & 禁止外推范围 \\
    \midrule
    200 bit分块BCH & 查表 & 整帧 & R15主、D19备选 & 285 & 明显2\%/5\%连续突发 & 不扩展到整块或300 bit BCH \\
    整块BCH & BM+Chien & 整帧 & NONE & 0 & 常规分散错误 & 无正式交织结论 \\
    CC常规 & 浮点软Viterbi，硬判决备选 & 滑窗/$50\times6$ & NONE & 0 & 常规信道与连续输出 & 参数按具体码率选择 \\
    CC 2\%突发 & 软Viterbi & 整帧 & D8/缓存优先 & 128 & 64 trellis-step跨度 & 与Pseudo128非等资源比较 \\
    CC 2\%突发 & 软Viterbi & 整帧 & Pseudo128/FER优先 & 256 & 128 trellis-step跨度 & 5\%下无稳定低FER恢复 \\
    LDPC & layered NMS，BP参考 & 码块迭代 & NONE & 0 & 已验证无突发及第7章信道 & 第9章不支持突发性能推断 \\
    \bottomrule
  \end{tabular}}
\end{table}

\section{最终推荐方案}

最终配置遵循“主推荐、备选、条件切换”三级表达。低速业务先按payload和复杂度、可靠性或突发目标选择BCH；高速业务先按目标码率组与信道场景选择CC或LDPC，再叠加译码、连续组织与交织配置。高速吞吐优先可采用CC约$3/4$与$W=160,S=25,D=126$，但必须先验证业务工作点保有足够FER余量。未知强突发没有为了填满配置表而指定未经验证的赢家。

工程控制面应把payload长度、目标FER、目标码率组、信道状态类别、实时预算和突发状态作为显式输入，不以单一“最佳编码”标志代替完整配置。一次模式选择至少同时下发编码方案标识、发送长度、译码器类型、软信息格式、滑窗参数、交织模式和缓存需求；接收端只有在这些字段与发送端一致时才启动对应译码流程。BCH的分组/整块切换会改变码字组织与译码器，卷积码的打孔切换会改变去打孔图样和发送长度，LDPC码长切换会改变提升因子、校验图与$\alpha$，因此均不能只修改名义码率而沿用旧状态。

条件切换也应设置滞回和最小保持时间，避免信道估计在门限附近波动时频繁改变码型。若可靠性监测显示主方案不再覆盖目标，应先降码率或切换至对应可靠性候选；若链路余量恢复且发送占用成为主要约束，再转回效率候选。突发告警只有在受损模式与第9章验证边界一致时才启用交织；未知强突发则进入降级状态，不以未经验证的交织配置维持正常业务声明。

每次配置迁移需重新执行无噪声编译码一致性、实际发送长度、FER工作点、输出时序和缓存上限五项验收。对卷积码还应验证尾比特和连续状态边界，对LDPC验证综合校验与最大迭代退出，对交织链路验证交织/解交织互逆。该验收集合使表\ref{tab:final-recommendations}成为可执行配置入口，同时保留前文各项Formal结果的适用条件。

\begin{figure}[htbp]
  \centering
  \fbox{\begin{minipage}[c][0.22\textheight][c]{0.88\textwidth}
    \centering
    \textbf{待插入：Visio综合选型流程图}\\[0.6em]
    电文业务 $\rightarrow$ 低速BCH / 高速目标码率\\
    低速：200/300 bit $\rightarrow$ 复杂度、可靠性或突发目标\\
    高速：近$2/3$ / 近$1/2$ $\rightarrow$ 信道、实时性与吞吐目标\\
    输出：编码、译码、连续组织、交织及条件切换
  \end{minipage}}
  \caption{低速与高速电文信道编码综合选型流程}
  \label{fig:final-selection-visio}
\end{figure}

\begin{table}[htbp]
  \centering\tiny
  \caption{最终方案总表}
  \label{tab:final-recommendations}
  \resizebox{\textwidth}{!}{\begin{tabular}{c c l l l l l l l}
    \toprule
    业务 & payload & 主要目标/场景 & 主推荐编码 & 码率/长度 & 译码与连续组织 & 交织 & 备选 & 适用边界 \\
    \midrule
    低速 & 200 & 低复杂度 & $19\times$BCH(15,11,1) & 0.701754/285 & 查表/整帧 & NONE & -- & 分散错误 \\
    低速 & 200 & 常规可靠性 & BCH(511,385) & 0.613497/326 & BM+Chien/整帧 & NONE & -- & 发送与计算开销增加 \\
    低速 & 200 & 明显连续突发 & 分块BCH & 0.701754/285 & 查表/整帧 & R15 & D19 & 仅限已验证分块码型 \\
    低速 & 300 & 综合平衡 & BCH(511,421) & 0.769231/390 & BM+Chien/整帧 & NONE & -- & 无正式交织配置 \\
    低速 & 300 & 可靠性优先 & BCH(511,385) & 0.704225/426 & BM+Chien/整帧 & NONE & BCH(511,421) & 更长发送占用 \\
    高速 & 300 & 近$2/3$常规信道 & CC R2/3 & 0.6536/459 & 软Viterbi，$W128/S25/D84$ & NONE & N480 & 近码率比较 \\
    高速 & 300 & 近$2/3$已知遮挡 & LDPC N480 & 0.6250/480 & NMS，$\alpha=0.95$ & NONE & -- & 受损位置已知 \\
    高速 & 300 & 近$1/2$可靠性 & LDPC N640 & 0.46875/640 & NMS，$\alpha=0.80$ & NONE & CC R1/2 & 迭代与图规模较大 \\
    高速 & 300 & 连续处理/等待可预测 & CC R1/2 & 0.4902/612 & 软Viterbi，$W96/S16/D70$，$50\times6$ & NONE & N640 & 可靠性边界需复核 \\
    高速 & 300 & 吞吐优先 & CC R3/4 & 0.7353/408 & 软Viterbi，$W160/S25/D126$ & NONE & CC R2/3 & 仅在FER余量允许时 \\
    高速 & 300 & 2\%突发 & CC按既定码率 & 对应码率/长度 & 软Viterbi & Pseudo128或D8 & -- & FER优先或缓存优先 \\
    高速 & 300 & 未知5\%强突发 & 无已验证稳定低FER主方案 & -- & -- & -- & -- & 需增强链路机制 \\
    \bottomrule
  \end{tabular}}
\end{table}

\section{全文主要结论、适用边界与后续建议}

\subsection{全文主要结论}

本项目形成以下系统结论。

第一，低速与高速业务不能使用统一编码方案。低速电文更适合按固定短消息组织BCH，高速电文则需要在连续处理卷积码和短块LDPC之间按目标码率与信道条件切换。

第二，BCH形成两条互补路线：短分组方案以查表译码获得较低实现复杂度，较长整块缩短方案以BM和Chien搜索获得更高常规信道可靠性。对已验证的200 bit分块方案，规则结构交织能显著改善集中突发错误，但该结论不属于所有BCH结构的固有属性。

第三，卷积码的核心工程价值是固定64状态、母码与打孔共用网格、可滑窗处理及按时隙分批输出。软判决应作为性能主路线，Q8可作为当前定点实现优先配置；硬判决只承担资源或接口受限时的备选角色。

第四，LDPC N640以较低实际码率和更大校验图换取更完整的低FER覆盖，适合近$1/2$可靠性和弱信号优先场景；其迭代次数、P95和软件处理代价也必须纳入系统预算。N480侧重传输效率，N560主要承担特定长度接口适配。

第五，NMS成为当前LDPC工程主译码器的原因，是在接近BP可靠性的同时显著降低当前软件处理时间；BP仍是必要的性能参考。该选择不是跨平台不变的算法优劣结论。

第六，高速近$2/3$和近$1/2$不能统一采用同一码族。常规全帧信道的近$2/3$组偏向CC R2/3，已知短时遮挡可切换N480；近$1/2$可靠性优先偏向N640，连续输出与可预测等待优先可选择CC R1/2。

第七，交织改变的是错误在码字中的空间分布，不改变码本及其冗余。200 bit分块BCH在已验证突发条件下收益明确，卷积码在2\%突发下可按FER或缓存选择配置，但5\%条件下尚未形成稳定低FER方案；LDPC没有建立交织或突发性能结论。

第八，未知强突发和严重连续损伤不能仅依靠增加FEC冗余解决。当前无目标覆盖的场景需要结合受损位置估计、擦除信息、重传或更有针对性的重排机制，而不是从现有常规信道结果直接指定主方案。

\subsection{当前适用边界}

当前结论主要建立在离散符号级BPSK链路上，尚未完整建模真实采样率、过采样、脉冲成形、射频前端、实际定时同步与载波跟踪、真实导航信号波形和硬件流水线。BCH多径结论依赖第4章规定的固定多径与MMSE接收条件；频偏结论依赖第7章规定的帧内相位轨迹，不能直接覆盖任意动态频偏过程。

LDPC方案是300 bit payload上的BG2 Direct QC短块实现，不等同于完整5G NR物理层链路。当前链路未纳入完整CRC、rate matching、bit selection及HARQ流程。BP和NMS的CPU时间只反映当前处理器、编译器、数据结构和计时范围；BCH、CC与LDPC操作语义不同，不能由软件数值推导硬件上必然存在的固定倍数关系。

交织缓存以coded bits计数。在没有冻结真实coded-bit rate、总线宽度和流水线并行度之前，285、128或256 bit缓存不能直接换算为物理毫秒。第9章的突发结论仅覆盖指定码型、2\%与5\%连续异常、规定位置集合及正式信噪比网格。5\%未知强突发中部分方案没有目标覆盖，结论不能外推到更长、更强或不同统计结构的突发过程。

\subsection{后续实现与验证建议}

\enlargethispage{5\baselineskip}

后续工作应优先把表\ref{tab:final-recommendations}中的低速BCH、高速CC和高速LDPC主方案接入统一端到端业务链，以相同电文接口、帧管理和接收状态输出验证条件切换。第二步应引入真实符号率与采样率，把coded-bit buffer、时隙等待和软件处理时间共同换算为系统可核验的物理时延预算。

第三步完善波形链路，包括脉冲成形、匹配滤波、定时同步、载波同步和信道估计，并在一致接收机条件下复核固定多径与频偏结论。第四步面向FPGA、DSP或SoC实现重新测量吞吐、逻辑与存储资源、端到端处理时延和功耗，分别验证查表/BM+Chien、Viterbi和layered NMS的数据通路。

第五步针对未知连续突发和强遮挡研究受损位置估计、擦除辅助、规则交织及必要的重传机制，先建立稳定目标覆盖再讨论效率优化。若系统要求标准兼容，第六步再补齐LDPC的CRC、rate matching、bit selection及HARQ相关链路，使当前Direct短块结果能够在明确接口下迁移到完整标准流程。
